% ASSERT_CONVENTION: natural_units=natural, metric_signature=mostly_minus, fourier_convention=physics, coupling_convention=alpha_s, generator_normalization=T(fund)=1/2, covariant_derivative_sign=D_mu=partial_mu-igA
%
% Lecture 11: The Koide Relation and the Dualised Standard Model
% Part of: The Dualised Standard Model -- Part III Lecture Notes
%
\documentclass[11pt,a4paper]{article}
% ASSERT_CONVENTION: natural_units=natural, metric_signature=mostly_minus, fourier_convention=physics, coupling_convention=alpha_s, generator_normalization=T(fund)=1/2, covariant_derivative_sign=D_mu=partial_mu-igA
%
% CONVENTIONS (Seiberg hep-th/9411149)
% Metric: (+,-,-,-) (mostly minus)
% T(fund) = 1/2 per Weyl fermion
% b_0 = (11N_c - 2N_f)/3 for N_f Dirac flavours
% D_mu = partial_mu - ig A_mu^a T^a
% A(fund) = 1 (cubic anomaly coefficient per fundamental Weyl fermion)
% alpha_s = g_s^2 / (4 pi)
% R[Q] = (N_f - N_c)/N_f; R[W] = 2
% N_f counts Dirac flavour pairs (each = 2 Weyl fermions)
%
% This file is \input{} by all lecture files.
% Lecture-local convention boxes and header blocks (e.g. Lectures 9-10) extend
% this shared base: they add lecture-specific conventions and keep every entry above.

%% ---------- Document class ----------
% (loaded by the wrapper; preamble only provides packages and macros)

%% ---------- Standard packages ----------
\usepackage{amsmath,amssymb,amsthm,mathtools}
\usepackage{hyperref}
\usepackage[capitalise,noabbrev]{cleveref}
\usepackage{enumitem}
\usepackage{booktabs}
\usepackage{array}
\usepackage{caption}                % \captionof for tables outside floats
\usepackage{xcolor}
\usepackage{tikz}
\usetikzlibrary{arrows.meta,decorations.markings,positioning,calc}
\usepackage{fancybox}
\usepackage{tcolorbox}
\tcbuselibrary{skins,breakable}
\usepackage{slashed}              % Feynman slash notation
\usepackage{cite}                 % compressed citation lists

%% ---------- Hyperref configuration ----------
\hypersetup{
  colorlinks=true,
  linkcolor=blue!70!black,
  citecolor=green!50!black,
  urlcolor=blue!80!black,
  pdfauthor={Part III Lecture Notes},
  pdftitle={The Dualised Standard Model}
}

%% ---------- Theorem environments ----------
\theoremstyle{plain}
\newtheorem{theorem}{Theorem}[section]
\newtheorem{lemma}[theorem]{Lemma}
\newtheorem{proposition}[theorem]{Proposition}
\newtheorem{corollary}[theorem]{Corollary}

\theoremstyle{definition}
\newtheorem{definition}[theorem]{Definition}
\newtheorem{example}[theorem]{Example}
\newtheorem{exercise}{Exercise}[section]

\theoremstyle{remark}
\newtheorem{remark}[theorem]{Remark}

%% ---------- Physics macros: gauge groups and representations ----------
\newcommand{\Nc}{N_{\!c}}
\newcommand{\Nf}{N_{\!f}}
\newcommand{\SU}[1]{\mathrm{SU}(#1)}
\newcommand{\UU}[1]{\mathrm{U}(#1)}

% Representations
\newcommand{\fund}{\square}
\newcommand{\antifund}{\overline{\square}}
\newcommand{\adj}{\mathrm{adj}}
\newcommand{\antisym}{\wedge^{\!2}}        % antisymmetric rank-2
\newcommand{\sym}{\mathrm{S}_2}           % symmetric rank-2

%% ---------- Physics macros: group theory constants ----------
% Dynkin index: Tr(T^a T^b) = T(R) delta^{ab}
\newcommand{\Tfund}{\tfrac{1}{2}}          % T(fund) = 1/2 per Weyl fermion
\newcommand{\Tadj}{\Nc}                    % T(adj) = N_c
\newcommand{\Cfund}{\frac{\Nc^2 - 1}{2\Nc}}  % C_2(fund)
\newcommand{\Cadj}{\Nc}                    % C_2(adj)

% Cubic anomaly coefficient: A(R) = Tr_R[T^a {T^b, T^c}] with A(fund) = 1
\newcommand{\Afund}{1}

%% ---------- Physics macros: beta function ----------
\newcommand{\bzero}{b_0}
\newcommand{\bone}{b_1}
\newcommand{\alphas}{\alpha_s}
\newcommand{\gs}{g_s}

%% ---------- Physics macros: SUSY / R-charge ----------
\newcommand{\Rcharge}[1]{R[#1]}
\newcommand{\Wsup}{W}                      % superpotential

%% ---------- Physics macros: covariant derivative and field strength ----------
\newcommand{\Dmu}{D_\mu}
\newcommand{\Fmn}{F_{\mu\nu}}
\newcommand{\Ftilde}{\widetilde{F}_{\mu\nu}}

%% ---------- Physics macros: common abbreviations ----------
\newcommand{\ie}{\textit{i.e.}}
\newcommand{\eg}{\textit{e.g.}}
\newcommand{\cf}{\textit{cf.}}
\newcommand{\IR}{\ensuremath{\mathrm{IR}}}
\newcommand{\UV}{\ensuremath{\mathrm{UV}}}
\newcommand{\AF}{\ensuremath{\mathrm{AF}}}
\newcommand{\BZ}{\ensuremath{\mathrm{BZ}}}
\newcommand{\NSVZ}{\ensuremath{\mathrm{NSVZ}}}
\newcommand{\SQCD}{\ensuremath{\mathrm{SQCD}}}
\newcommand{\QCD}{\ensuremath{\mathrm{QCD}}}
\newcommand{\SM}{\ensuremath{\mathrm{SM}}}
\newcommand{\Tr}{\mathrm{Tr}}

%% ---------- Convention box macro ----------
\newtcolorbox{conventionbox}[1][]{%
  enhanced,
  colback=blue!5!white,
  colframe=blue!60!black,
  fonttitle=\bfseries,
  title={Conventions},
  #1
}

%% ---------- Comparison table style ----------
\newtcolorbox{comparisontable}[1][]{%
  enhanced,
  colback=orange!5!white,
  colframe=orange!70!black,
  fonttitle=\bfseries,
  title={Comparison},
  #1
}

%% ---------- Warning / important box ----------
\newtcolorbox{warningbox}[1][]{%
  enhanced,
  colback=red!5!white,
  colframe=red!70!black,
  fonttitle=\bfseries,
  title={Important},
  #1
}

%% ---------- Preview / forward-reference box ----------
\newtcolorbox{previewbox}[1][]{%
  enhanced,
  colback=green!5!white,
  colframe=green!50!black,
  fonttitle=\bfseries,
  title={Preview},
  #1
}

%% ---------- Page layout ----------
\usepackage[margin=2.5cm]{geometry}
\setlength{\parskip}{0.4em}
\setlength{\parindent}{0em}

%% ---------- Section formatting ----------
\usepackage{titlesec}
\titleformat{\section}{\Large\bfseries}{\thesection}{1em}{}
\titleformat{\subsection}{\large\bfseries}{\thesubsection}{1em}{}

%% ---------- End of preamble ----------


\title{\textbf{Lecture 11: The Koide Relation and the Dualised
Standard Model}}
\author{The Dualised Standard Model --- Part III Lecture Notes}
\date{Lent Term 2027}

\begin{document}
\maketitle

\tableofcontents
\newpage

%% ========================================================================
\section{Introduction and motivation}
\label{sec:intro}
%% ========================================================================

In Lectures~7--8, we constructed the dualised Standard Model: the
SM, colour group $\SU{3}_c$ included, as the conjectured
\emph{magnetic} dual of an electric $\SU{3}$ gauge theory (a group
distinct from colour) with $\Nf = 6$
flavours~\cite{Sannino2024}.  Eighteen Higgs doublets emerge from
the mes-Higgs field $\Phi_H$, and the fermion mass hierarchy arises
from hierarchical VEVs among these doublets.  The Yukawa coupling
$y$ is \emph{universal} (fixed by the magnetic Lagrangian of
Lecture~7), and the effective Yukawa matrix is
\begin{equation}
  y_{ij}^{\text{eff}} = \sqrt{2}\,y\,\frac{\langle H_{ij}\rangle}{v}\,,
  \label{eq:yeff}
\end{equation}
where $H_{ij}$ are the Higgs doublet components, $v = 246$~GeV is
the electroweak VEV, and the SM Higgs $H_0 \equiv H^u_{33}$ carries
$\langle H_0\rangle \simeq v/\sqrt{2}$, so that $y^{\text{eff}}_{33}
\simeq y$.

In this lecture we ask: \emph{does the dualised SM predict or
constrain the fermion mass eigenvalues?}  We establish one solid
result and explore two conjectures:
\begin{enumerate}
\item[\textbf{(R)}] The Koide relation $Q = 3/2$ is algebraically
  equivalent to a \textbf{VEV condition} on the Higgs
  doublet VEV ratios---a reformulation that introduces no new
  parameters and follows from the universal Yukawa structure of
  Lecture~8.
\item[\textbf{(C1)}] The \emph{numerical value} $Q = 3/2$ coincides
  with the SU(5) Casimir ratio
  $C_2(\mathbf{10})/C_2(\mathbf{5}) = 2(N-2)/(N-1)$, which
  uniquely selects $N=5$.  Whether the Higgs potential enforces this
  value is \emph{conjectural}.
\item[\textbf{(C2)}] The split of the five light quarks into
  $k_u = 2$ up-type and $k_d = 3$ down-type flavours, together with
  $N_{\text{gen}} = 3$, solves the sBootstrap closure conditions.
  The ISS rank condition reproduces the same integers at
  $(\Nc,\Nf) = (3,5)$; in the duality frame of Lecture~7 this match
  is numerical and supplies no mechanism (\S\ref{sec:sbootstrap}).
\end{enumerate}
The Koide angle $\delta = 2/9$ remains an empirically determined
parameter; its origin is an open problem (see~\S\ref{sec:angle}).

%% ========================================================================
\section{The Koide relation as a VEV condition}
\label{sec:koide-vev}
%% ========================================================================

This section contains the lecture's motivating reformulation.  It requires only
the universal Yukawa coupling of Lecture~8 and standard algebra.
The observation that the Koide condition becomes a VEV constraint
in universal-Yukawa models was first made in the yukawaon
programme~\cite{Koide1983,HabaKoide2008,Ma2007}; we present it
here in the specific context of the dualised SM.

\subsection{Setup: from VEVs to masses}

From Lecture~8, the $k$-th dormant doublet acquires the induced VEV
$\langle H_k\rangle = (\mu_{0k}^2/M_k^2)\,\langle H_0\rangle$, where
$\mu_{0k}^2$ is the mixing mass between the SM Higgs $H_0$ and the
$k$-th dormant doublet, and $M_k$ is the diagonal mass of the $k$-th
doublet.  With $\langle H_0\rangle = v/\sqrt{2}$, the effective
Yukawa coupling for the $k$-th generation in a given charge sector
is
\begin{equation}
  y_k^{\text{eff}} = \sqrt{2}\,y\,\frac{\langle H_k\rangle}{v}
  = y\,\frac{\mu_{0k}^2}{M_k^2}\,.
  \label{eq:yk}
\end{equation}
The fermion mass is then
\begin{equation}
  m_k = y_k^{\text{eff}}\,\frac{v}{\sqrt{2}}
  = y\,\frac{v}{\sqrt{2}}\,\frac{\mu_{0k}^2}{M_k^2}\,.
  \label{eq:mk}
\end{equation}

Define the \emph{mixing ratio}:
\begin{equation}
  x_k \equiv \frac{\mu_{0k}}{M_k}\,.
  \label{eq:xk}
\end{equation}
Then $m_k = m_0\,x_k^2$ where $m_0 = y\,v/\sqrt{2}$ is the
universal scale, and
\begin{equation}
  \sqrt{m_k} = \sqrt{m_0}\,x_k\,.
  \label{eq:sqrtm}
\end{equation}

\subsection{The Koide condition on mixing ratios}

The Koide parameter for three fermion masses $m_1, m_2, m_3$ is
\begin{equation}
  Q = \frac{(\sqrt{m_1}+\sqrt{m_2}+\sqrt{m_3})^2}
  {m_1+m_2+m_3}\,.
  \label{eq:Q}
\end{equation}
Substituting~\eqref{eq:sqrtm}:
\begin{equation}
  \boxed{Q = \frac{(x_1+x_2+x_3)^2}{x_1^2+x_2^2+x_3^2}\,.}
  \label{eq:Qx}
\end{equation}
The universal scale $m_0$ cancels.  The Koide condition
$Q = 3/2$ is a \textbf{pure constraint on the mixing ratios}
$x_k = \mu_{0k}/M_k$.  This is the key rewriting of this
lecture: in any model with a universal Yukawa coupling and
multiple Higgs doublets, the Koide relation reduces to a
condition on the \emph{Higgs sector} alone.  The cancellation
of $m_0$ follows from the scale invariance of $Q$ (homogeneous of
degree zero) and therefore holds in every such model, the dualised SM
included.

\begin{warningbox}
The Koide condition is independent of the universal Yukawa
coupling $y$ and the Higgs VEV $v$.  It constrains only the
\emph{ratios} $\mu_{0k}/M_k$ of the Higgs sector parameters.
This reformulation is exact.  What follows in Sections~3--4
is conjectural.
\end{warningbox}

\subsection{Decomposition into common and traceless parts}

Write $x_k = z_0 + z_k$ with
\begin{equation}
  z_0 = \frac{x_1+x_2+x_3}{3}\,,\qquad
  z_1+z_2+z_3 = 0\,.
  \label{eq:z0zk}
\end{equation}
Then:
\begin{equation}
  Q = \frac{9\,z_0^2}{3\,z_0^2 + \sum_k z_k^2}\,.
  \label{eq:Qz}
\end{equation}
The condition $Q = 3/2$ requires
\begin{equation}
  z_0^2 = \frac{1}{3}\sum_{k=1}^3 z_k^2\,.
  \label{eq:koide-cond}
\end{equation}
This is the \textbf{trace condition}: the squared common piece
equals one-third of the sum of squared traceless pieces.
It is identical to Koide's original 1983 constraint
$\Tr[U^2] = \Tr[V^2]$~\cite{Koide1983}.

\subsection{What the potential must do}

The effective potential for the dormant Higgs doublets
(Lecture~8) is
\begin{equation}
  V_{\text{eff}} = \sum_k M_k^2\,|H_k|^2
  - \sum_k \mu_{0k}^2\,(H_0^\dagger H_k + \text{h.c.})
  + \text{quartic terms}\,.
  \label{eq:Veff}
\end{equation}
Minimising with respect to $H_k$ at fixed $\langle H_0\rangle$ gives
$\langle H_k\rangle/\langle H_0\rangle = \mu_{0k}^2/M_k^2 = x_k^2$,
as used above.  The quartic self-couplings of the dormant doublets
shift this only at relative order $\langle H_k\rangle^2/M_k^2$, which
Lecture~8 neglects; quartic couplings to $H_0$ renormalise $M_k^2$
and $\mu_{0k}^2$ by amounts of order $\lambda v^2$.  The Koide
condition~\eqref{eq:koide-cond} is therefore a \textbf{constraint on
the (renormalised) mass parameters} $\mu_{0k}$ and $M_k$ of the
potential: whatever generates them
(Planck-suppressed operators of the electric theory, radiative
corrections, or quartic couplings to other condensates) must deliver
ratios $x_k$ with $z_0^2 = \frac{1}{3}\sum z_k^2$.

We do \emph{not} derive that the potential has this property.
Whether the emergent quasi-SUSY of the magnetic theory constrains
these parameters sufficiently is the central open problem of this
programme.  Exercise Set~9 (Problem~5) shows that an $S_3$-symmetric
potential with generic soft breaking leaves $Q$ free.

%% ========================================================================
\section{Conjecture C1: $Q = 3/2$ and the SU(5) Casimir ratio}
\label{sec:casimir}
%% ========================================================================

\subsection{The numerical coincidence}

For $\SU{N}$, the quadratic Casimir operators of the fundamental
representation~($\fund$) and the antisymmetric two-tensor
($\antisym\fund$) satisfy:
\begin{equation}
  \frac{C_2(\antisym\fund)}{C_2(\fund)}
  = \frac{2(N-2)}{N-1}\,.
  \label{eq:casimir-ratio}
\end{equation}
This follows from $C_2(\wedge^p\,N) = p(N-p)(N+1)/(2N)$.
Setting the ratio to $3/2$:
\begin{equation}
  \frac{2(N-2)}{N-1} = \frac{3}{2}
  \qquad\Longrightarrow\qquad N = 5\,.
  \label{eq:Q-N5}
\end{equation}
This is the unique positive integer solution.

\begin{center}\small
\begin{tabular}{@{}ccc@{}}
\toprule
$N$ & $C_2(\wedge^2 N)/C_2(N)$ & \\
\midrule
$3$ & $1$ & (degenerate) \\
$4$ & $4/3$ & \\
$\mathbf{5}$ & $\mathbf{3/2}$ & $= Q_{\text{Koide}}$ \\
$6$ & $8/5$ & \\
$\infty$ & $2$ & \\
\bottomrule
\end{tabular}
\end{center}

\subsection{Why $N = 5$?}

In the dualised SM of Lectures~7--8, the flavour group is
$\SU{6}_L \times \SU{6}_R$ (from $\Nf = 6$).  The value $N = 5$
points to an $\SU{5}$ subgroup of the flavour $\SU{6}$.
The five lightest quarks $(u,c,d,s,b)$ form the fundamental
of an SU(5) flavour subgroup, with the top quark as the sixth
flavour playing a special role.

This SU(5) flavour structure is precisely the \emph{sBootstrap}
framework~\cite{Rivero2009,Rivero2024}, where the top is called
the ``elephant'' and the five light quarks are ``turtles.''
The decomposition
\begin{equation}
  \SU{6}_{\text{flavour}} \supset \SU{5}_{\text{turtle}}
  \times \UU{1}_{\text{top}}
\end{equation}
is natural from the mass hierarchy alone ($m_t \gg m_b$), but
the duality does not \emph{derive} it: nothing in the CS framework
selects SU(5) over SU(6) as the relevant Casimir group.

\begin{warningbox}
The coincidence $C_2(\mathbf{10})/C_2(\mathbf{5}) = 3/2
= Q_{\text{Koide}}$ is a \emph{mathematical fact} about SU(5).
Whether the Higgs potential minimum is determined by this
Casimir ratio is a \emph{conjecture without proof}.  No
derivation exists connecting the one-loop potential of the
18-doublet sector to the eigenvalue condition~\eqref{eq:koide-cond}.
\end{warningbox}

For this coincidence to be physically meaningful, there must exist
a dynamical channel through which $C_2$ ratios in an $\SU{5}$
flavour group influence the Higgs potential minimum.  Candidate
mechanisms include: one-loop corrections proportional to $C_2$
entering the quartic couplings, renormalisation-group fixed-point
conditions, or representation-theoretic constraints on the
potential structure.  None of these has been investigated.

\subsection{Dynkin index integers}

If one takes the SU(5) Casimir connection seriously, the
Dynkin index ratios provide further numerical coincidences:
\begin{align}
  \frac{T(\mathbf{15})}{T(\mathbf{5})} &= N+2 = 7\,,
  \label{eq:dyn-7} \\
  \frac{T(\mathbf{24})}{T(\mathbf{5})} &= 2N = 10\,,
  \label{eq:dyn-10} \\
  \frac{T(\mathbf{10})}{T(\mathbf{5})} &= N-2 = 3 = \Nc\,.
  \label{eq:dyn-3}
\end{align}
The integers $7$ and $10$ also appear in the empirical
``$Z$-ladder'' of Ref.~\cite{RiveroZladder}: pairwise mass relations
$m_X^2/m_Y = M_Z/n^2$ with small integers $n$, for example
$m_s \simeq \sqrt{m_d M_Z}/7$ ($93.52$ against $93.5$~MeV) and
$m_c \simeq \sqrt{m_\tau M_Z}/10$ ($1272.9$ against $1273$~MeV),
built from masses at mixed scales.  The integer $N-2 = 3$ is the
number of QCD colours.  These coincidences are suggestive but do not
constitute a derivation.

%% ========================================================================
\section{The Koide angle $\delta$: an open problem}
\label{sec:angle}
%% ========================================================================

The Koide condition~\eqref{eq:koide-cond} constrains the sum
$\sum z_k^2$ relative to $z_0^2$, but does not fix the
\emph{individual} $z_k$ values.  The mass spectrum is
parametrised by a single angle $\delta$ in the $Z_3$
representation:
\begin{equation}
  x_k = z_0\!\left(1+\sqrt{2}\,
  \cos\!\left(\frac{2\pi k}{3}+\delta\right)\right),
  \qquad k = 1,2,3\,.
  \label{eq:z3}
\end{equation}
The value $\delta = 2/9$~rad reproduces the charged lepton
masses to $10^{-5}$ precision~\cite{Brannen2006}.

Neither the dualised SM nor the sBootstrap currently derives
this value.  Within the multi-Higgs framework, $\delta$ is
determined by the relative magnitudes of the mixing masses
$\mu_{0k}$ and the diagonal masses $M_k$.  These originate
from Planck-scale operators in the electric theory
(Lecture~8), and their structure is not sufficiently
constrained by the duality to fix $\delta$.

Deriving $\delta$ from a dynamical principle---whether from
modular flavour symmetries, conformal fixed-point dynamics,
or some other mechanism---remains the central open problem
of the Koide programme.

%% ========================================================================
\section{What $Q = 3/2$ determines and what it leaves free}
\label{sec:spectrum}
%% ========================================================================

The Koide condition $Q = 3/2$ is one constraint on three
eigenvalues, leaving a two-parameter family: the overall
scale $m_0$ and the angle $\delta$.  The scale $m_0$ sets
the heaviest mass; $\delta$ sets the ratios.

With the empirical $\delta = 2/9$, the Koide formula predicts
\begin{align}
  \frac{m_e}{m_\mu} &= 0.004836\,,
  & \text{(measured: } 0.004836\text{)} \\
  \frac{m_\mu}{m_\tau} &= 0.05946\,,
  & \text{(measured: } 0.05946\text{)}
\end{align}
They agree with the measured ratios to $1.0\times10^{-5}$ and
$2.1\times10^{-5}$ (relative).  This precision is the content of the
original Koide formula~\cite{Koide1983} with $\delta = 2/9$.  The
dualised SM so far provides $Q = 3/2$ as a VEV condition and leaves
$\delta$ free.

%% ========================================================================
\section{The chain angle and colour}
\label{sec:chain}
%% ========================================================================

The Koide chain~\cite{Rivero2013,Zenczykowski2012}
extends the lepton relation to quarks:
$(0,s,c) \to (-s,c,b) \to (c,b,t)$, with each link
satisfying $Q \approx 3/2$ and the chain angle
$\delta_{scb} \approx 3\,\delta_\ell$.

One interpretation of the factor~$3$: the quark sector carries
SU(3)$_c$ colour, which multiplies the vacuum angle:
\begin{equation}
  \delta_{\text{chain}} = \Nc\,\delta_\ell
  = 3\times\frac{2}{9} = \frac{2}{3}\;\text{rad}
  = 38.2^\circ\,.
  \label{eq:chain-angle}
\end{equation}
The measured value $\delta_{scb} \approx 37.2^\circ$ agrees
to $2.7\%$.  No calculation within the dualised SM or the
sBootstrap currently produces the factor $\Nc$ from dynamics.
Whether the $2.7\%$ agreement is evidence for the colour
interpretation or a numerical coincidence cannot be decided
without a derivation of the vacuum structure in the coloured
sector.

%% ========================================================================
\section{The inverse Koide from the seesaw}
\label{sec:inverse}
%% ========================================================================

\subsection{$Q$ is preserved under inversion}

In SQCD with every flavour massive, holomorphy fixes the meson VEVs
to be inversely proportional to the quark masses,
$\langle Q_k\widetilde Q^k\rangle = C/m_k$ with the flavour-independent
constant $C = (\det m\,\Lambda^{3\Nc-\Nf})^{1/\Nc}$~\cite{IntriligatorSeiberg1995}:
a seesaw between quark masses and meson VEVs.  Since $Q$ is
homogeneous of degree zero:
\begin{equation}
  Q\!\left(\frac{C}{m_1},\frac{C}{m_2},\frac{C}{m_3}\right)
  = Q\!\left(\frac{1}{m_1},\frac{1}{m_2},\frac{1}{m_3}\right).
  \label{eq:Q-inverse}
\end{equation}
This is an exact identity.

\subsection{Empirical status}

With PDG~2024 pole masses for the leptons and the
$\overline{\text{MS}}$ quark masses $m_d(2\;\text{GeV}) = 4.70$~MeV,
$m_s(2\;\text{GeV}) = 93.5$~MeV and $m_b(m_b) = 4183$~MeV~\cite{AHS2025}:
\begin{equation}
  Q(m_e,m_\mu,m_\tau) = 1.500005\,,\qquad
  Q(1/m_d,1/m_s,1/m_b) = 1.5045 \pm 0.010\,.
\end{equation}
The uncertainty on $Q(1/m)$ is dominated by $m_d = 4.70 \pm 0.20$~MeV
($4.3\%$).  The deviation from $3/2$ is $0.0045 \pm 0.010$, i.e.\
within $1\sigma$: the agreement cannot currently be distinguished
from numerical coincidence at this level of quark-mass precision.

Note that Q-preservation under inversion is a mathematical identity
holding for any mass triplet; the empirical near-agreement
$Q(1/d,1/s,1/b) \approx Q(e,\mu,\tau) \approx 3/2$ is a separate
observation requiring an independent dynamical explanation.

The \emph{individual mass ratios} do not match through a simple
seesaw $m_\ell = c/m_q$ (see~\S\ref{sec:caveats}).

\subsection{The hierarchy problem of the seesaw}
\label{sec:caveats}

The seesaw $m_\ell \propto 1/m_q$ gives an \emph{inverted}
hierarchy: the lightest quark maps to the heaviest lepton.
With the standard generation pairing $(d \leftrightarrow e,\;
s \leftrightarrow \mu,\; b \leftrightarrow \tau)$:
$m_e/m_\mu = m_s/m_d = 19.9$, while the actual ratio is
$0.0048$---a factor of $4000$ discrepancy.

With the reversed pairing $(d \leftrightarrow \tau,\;
s \leftrightarrow \mu,\; b \leftrightarrow e)$: $m_\tau$
is recovered by construction, $m_\mu$ is $15\%$ off, and
$m_e$ is a factor~$4$ wrong.

The seesaw preserves $Q$ (that is exact) but does \emph{not}
reproduce the individual lepton masses.  The quark--lepton
mass mapping is more complex than a simple inversion.

%% ========================================================================
\section{Connection to the sBootstrap}
\label{sec:sbootstrap}
%% ========================================================================

The sBootstrap framework~\cite{Rivero2009,Rivero2024} identifies
sfermions with known mesons and diquarks via $\SU{5}$ flavour
symmetry.  The closure conditions~\cite{Rivero2024}
\begin{equation}
  2N = k_u\,k_d\,,\quad
  2N = \frac{k_d(k_d+1)}{2}\,,\quad
  4N = k_u^2+k_d^2-1
\end{equation}
have the unique positive solution $N = 3$ generations with
$k_u = 2$, $k_d = 3$ (Exercise Set~9, Problem~4).

The ISS rank condition~\cite{ISS2006} reproduces the same two
integers for SQCD with $(\Nc,\Nf) = (3,5)$.  The next subsection
places that counting in the duality frame of Lecture~7 and shows
that the match stays numerical.

\subsection{From $\Nf = 6$ to $\Nf = 5$: decoupling the top quark}
\label{sec:top-decoupling}

\paragraph{The frame.}  In Lecture~7 (\S1 and \S3) the electric
theory is an $\SU{3}$ gauge theory, distinct from colour, with
$\Nf = 6$ flavours $Q$, $\widetilde Q$.  Its conjectured magnetic dual
has gauge group $\SU{\Nf-\Nc} = \SU{3}_c$, magnetic quarks $q$,
$\widetilde q$ identified with the SM quarks, and a $6\times 6$ meson
$M$ with 36~components whose scalar partner $\Phi_H$ contains the
Higgs doublets.  In the Pati--Salam route of Lecture~7 (\S4) the
electric group is $\SU{2}$ and the magnetic group is
$\SU{4}_{\text{PS}} \supset \SU{3}_c$.  Colour is magnetic in both
routes.

\paragraph{The ISS rank condition.}  If every electric flavour
carries a small mass $m$, the magnetic superpotential acquires a
linear term~\cite{ISS2006},
\begin{equation}
  W_{\text{mag}} = h\,\Tr(q\,M\,\widetilde q) - h\,\mu^2\,\Tr M\,,
  \qquad \mu^2 \sim m\,\Lambda\,.
  \label{eq:iss-w}
\end{equation}
The F-term conditions $\widetilde q\,q = \mu^2\,\mathbb{1}_{\Nf}$
cannot all hold: the $\Nf\times\Nf$ matrix $\widetilde q\,q$ has rank
at most $\Nf-\Nc$, the number of magnetic colours.  In the
metastable vacuum $\widetilde q\,q$ condenses along $\Nf-\Nc$ flavour
directions, which Higgses the magnetic gauge group completely, and
the remaining $\Nc$ directions carry the nonzero F-terms that break
supersymmetry.  This is the $(\Nf-\Nc)+\Nc$ split.  ISS establish the
vacuum in the free-magnetic range $\Nc < \Nf < \tfrac{3}{2}\Nc$,
where the magnetic theory is infrared free.

\paragraph{Two ways to remove the top.}  The two deformations that
take $\Nf = 6$ to five light flavours act differently
(Lecture~3, \S5; \cite{Seiberg1995}):
\begin{itemize}
\item An electric mass $m_t\,Q_t\widetilde Q_t$ maps to a linear term
  $m_t\,M_{tt}$ in $W_{\text{mag}}$.  Its F-term forces
  $\langle q_t\,\widetilde q_t\rangle \neq 0$, and this dual-quark
  VEV Higgses the magnetic group:
  \begin{equation}
    \SU{3}_c \;\xrightarrow{\;\langle q_t\widetilde q_t\rangle \neq 0\;}
    \;\SU{2}\,.
    \label{eq:higgsing}
  \end{equation}
  In the dualised SM this breaks colour.
\item A meson VEV $\langle M_{tt}\rangle \neq 0$ gives the magnetic
  top quark a mass through the coupling $M q\widetilde q$ and leaves
  $\SU{3}_c$ intact.  This is how Lecture~7 generates $m_t$, through
  $\langle H^u_{33}\rangle$.  On the electric side it corresponds to
  $\langle Q_t\widetilde Q_t\rangle \neq 0$, which Higgses the electric
  $\SU{3} \to \SU{2}$; the dual of $\SU{2}$ with five flavours is
  again $\SU{5-2} = \SU{3}_c$.
\end{itemize}

\paragraph{What survives of the $2+3$ split.}  A tempting argument
takes $\SU{3}_c$ as the \emph{electric} group, gives the top an
electric mass, lets $\langle q_t\widetilde q_t\rangle$ Higgs a
magnetic $\SU{3}$ (then distinct from colour) down to $\SU{2}$, and
applies the rank condition to electric $\SU{3}_c$ with five
flavours:
\begin{equation}
  k_u = \Nf - \Nc = 2\,,\qquad k_d = \Nc = 3\,.
  \label{eq:iss-split}
\end{equation}
In the frame of Lecture~7 this chain fails at three points.
(i)~The electric top mass would break colour.
(ii)~The ISS condensate $\langle\widetilde q\,q\rangle$ Higgses the
magnetic group, which is again colour.
(iii)~$(\Nc,\Nf) = (3,5)$ lies inside the conformal window
$\tfrac{3}{2}\Nc < \Nf < 3\Nc$, outside the range where the ISS
vacuum is under control.  For the pair that remains after the top
decouples through $\langle H^u_{33}\rangle$ (electric $\SU{2}$ with
five flavours, magnetic $\SU{3}_c$) the same counting gives
$(\Nf-\Nc,\Nc) = (3,2)$, which fails the second closure condition.
The integers~\eqref{eq:iss-split} therefore match the sBootstrap
values numerically at $(\Nc,\Nf) = (3,5)$, and the dualised SM does
not derive them.

\begin{warningbox}
Status of $k_u = 2$, $k_d = 3$: \emph{conjectural}.  They solve the
sBootstrap closure conditions and coincide with the ISS counting
$\Nf-\Nc$, $\Nc$ at $(\Nc,\Nf) = (3,5)$.  A dynamical derivation
inside the dualised SM would need a mechanism that splits the five
light flavours without breaking $\SU{3}_c$; none is known.
\end{warningbox}

\subsection{The meson decomposition}

The $5\times 5$ meson matrix transforms as
$\mathbf{5}\otimes\overline{\mathbf{5}} = \mathbf{24}\oplus\mathbf{1}$
under $\SU{5}$ flavour.  The up/down split $(u,c\,|\,d,s,b)$
decomposes $\SU{5} \to \SU{2}_u \times \SU{3}_d \times \UU{1}$, and
the 24~adjoint mesons organise by electric charge:
\begin{center}\small
\begin{tabular}{@{}llcc@{}}
\toprule
Block & $\SU{2}\times\SU{3}$ rep & Charge & Count \\
\midrule
$\bar u\,u$\;(traceless) & $(\mathbf{3},\mathbf{1})$ & $0$ & 3 \\
$\bar d\,d$\;(traceless) & $(\mathbf{1},\mathbf{8})$ & $0$ & 8 \\
$\UU{1}$ generator & $(\mathbf{1},\mathbf{1})$ & $0$ & 1 \\
$\bar d\,u$ & $(\mathbf{2},\overline{\mathbf{3}})$ & $+1$ & 6 \\
$\bar u\,d$ & $(\overline{\mathbf{2}},\mathbf{3})$ & $-1$ & 6 \\
\bottomrule
\end{tabular}
\end{center}
giving $\mathbf{24} \to 12_0 + 6_{+1} + 6_{-1}$, plus the
flavour-singlet $\eta'$-like trace for a total of~25.

The sBootstrap~\cite{Rivero2009,Rivero2024} identifies these 25
composites with the scalar partners of known fermions: the charged
mesons ($\bar d_i u_j$ and $\bar u_i d_j$) map to charged
sfermions, while the neutral composites map to sneutrinos and
neutral Higgs-like states.  The closure conditions then
count:
\begin{equation}
  2N = k_u\,k_d = 6\,,\quad
  2N = \tfrac{k_d(k_d+1)}{2} = 6\,,\quad
  4N = k_u^2+k_d^2-1 = 12
\end{equation}
all giving $N = 3$ generations \emph{uniquely}.

\subsection{Koide from the meson mass formula}

The sBootstrap mass formula $m_k = (z_0+z_k)^2\,K_\Omega$
is algebraically equivalent to the Koide condition~\eqref{eq:koide-cond},
as shown in~\cite{Rivero2024}.  The dualised SM provides the
\emph{physical context} for this formula: $z_0+z_k$ is the
mixing ratio $x_k$ of the Higgs sector, and $K_\Omega$ is
the universal scale $m_0$.

%% ========================================================================
\section{Summary and honest assessment}
\label{sec:summary}
%% ========================================================================

\begin{center}
\renewcommand{\arraystretch}{1.3}
\begin{tabular}{lll}
\toprule
\textbf{Result} & \textbf{Origin} & \textbf{Status} \\
\midrule
Koide $=$ VEV condition & Universal Yukawa $+$ algebra
  & \textbf{Reformulated} \\
$Q(C/m_k) = Q(1/m_k)$ & Homogeneity of $Q$
  & \textbf{Theorem} \\
$Q(1/d,1/s,1/b) \approx 3/2$ & PDG masses, mixed scales
  & Within $1\sigma$ \\
\midrule
$k_u = 2$, $k_d = 3$ & sBootstrap closure; ISS count at $(3,5)$
  & Numerical match \\
$N_{\text{gen}} = 3$ & sBootstrap closure conditions
  & Counting (conjectural) \\
$Q = C_2(\mathbf{10})/C_2(\mathbf{5})$ & SU(5) Casimir
  & Observation \\
$\delta_{scb} = \Nc\,\delta_\ell$ & Empirical (factor 3)
  & $2.7\%$ match \\
$\delta = 2/9$ & --- & Open problem \\
\midrule
VEV condition mechanism & Origin of $\mu_{0k}$, $M_k$
  & \emph{Not derived} \\
Quark--lepton mapping & Seesaw details
  & \emph{Fails quantitatively} \\
$Z$-ladder relations & Empirical
  & \emph{Not derived} \\
\bottomrule
\end{tabular}
\end{center}

\textbf{Scale convention.}  The VEV condition is imposed at
the scale where the dormant Higgs doublets are integrated out,
their mass $M_k \sim \Lambda_S \sim$ few~TeV, where the
quasi-supersymmetric magnetic spectrum sets in (Lecture~7, \S4).
This scale lies far below the duality scale
$\Lambda_D \sim 10^{11}$~GeV at which the electric theory takes
over.  Below $\Lambda_S$, RG running introduces corrections that
spoil $Q = 3/2$ for quark masses at low energies.  The
charged-lepton Koide relation holds for \emph{pole masses}; the
quark chain uses $\overline{\text{MS}}$ running masses at mixed
scales.  This difference is a known subtlety~\cite{Sumino2009} and
is unresolved in the present framework.

\textbf{What is established.}  The universal Yukawa coupling of
the dualised SM translates the Koide relation into a VEV
condition on the multi-Higgs sector.  The seesaw between quark
masses and meson VEVs preserves $Q$ as a mathematical identity.

\textbf{What is observed but not derived.}  The numerical
coincidence $Q = C_2(\mathbf{10})/C_2(\mathbf{5}) = 3/2$ is an
exact group-theoretic fact that matches the Koide value.
It has no known dynamical origin within the dualised SM.
The turtle/elephant split $k_u = 2$, $k_d = 3$ matches the ISS
counting at $(\Nc,\Nf) = (3,5)$ only numerically
(\S\ref{sec:top-decoupling}).
The angle $\delta = 2/9$ has no known derivation at all.

\textbf{What remains open.}  The dynamics that fixes the ratios
$\mu_{0k}/M_k$ of the 18-Higgs potential must be shown to land on
the Koide point; the quartic couplings of the dormant doublets enter
the minimum only at relative order $\langle H_k\rangle^2/M_k^2$.
The quark--lepton mass mapping through the seesaw preserves $Q$ but
not the individual mass ratios.

%% ========================================================================
\section{Exercises}
\label{sec:exercises}
%% ========================================================================

Worked solutions to these five exercises are in Exercise Set~9.

\begin{exercise}
Verify that $C_2(\wedge^2\mathbf{N})/C_2(\mathbf{N})
= 2(N-2)/(N-1)$ using the general formula
$C_2(\wedge^p N) = p(N-p)(N+1)/(2N)$.  For which values
of $N$ is this ratio an integer?
\end{exercise}

\begin{exercise}
Compute $Q$ for the three charged lepton pole masses
$(m_e = 0.511\;\text{MeV},\; m_\mu = 105.658\;\text{MeV},\;
m_\tau = 1776.93\;\text{MeV})$.  Compare with
$Q(1/m_d,\,1/m_s,\,1/m_b)$ using the PDG values
$m_d = 4.70$, $m_s = 93.5$, $m_b = 4183\;\text{MeV}$.
\end{exercise}

\begin{exercise}
In the $Z_3$ parametrisation~\eqref{eq:z3}, show that
$Q = 3/2$ is satisfied for \emph{any} value of $\delta$.
Then verify that $\delta = 2/9$ reproduces the lepton mass
ratios $m_e/m_\mu$ and $m_\mu/m_\tau$ to better than
$0.01\%$.
\end{exercise}

\begin{exercise}
Show that the sBootstrap closure conditions give $k_u = 2$,
$k_d = 3$ uniquely.  The ISS rank condition splits the flavours of
SQCD with $\Nf$ flavours and $\Nc$ colours as $\Nf - \Nc$ and
$\Nc$.  For which values of $(\Nc, \Nf)$ do the two countings agree?
Does that pair lie in the range $\Nc < \Nf < \tfrac{3}{2}\Nc$ where
the ISS vacuum is under control, and which gauge group would the ISS
condensate Higgs in the dualised SM of Lecture~7?
\end{exercise}

\begin{exercise}[Challenge]
Write the most general quartic potential for three complex
scalar doublets $H_1, H_2, H_3$ with an $S_3$ permutation
symmetry softly broken by the mixing masses $\mu_{0k}$.
Does the minimum satisfy~\eqref{eq:koide-cond} for any
choice of quartic couplings?  What additional conditions
(if any) are needed?
\end{exercise}

%% ========================================================================
%% BIBLIOGRAPHY
%% ========================================================================
\begin{thebibliography}{20}

\bibitem{Sannino2024}
G.~Cacciapaglia and F.~Sannino,
``Charting Standard Model Duality and its Signatures,''
arXiv:2407.17281 [hep-ph] (2024).

\bibitem{Rivero2009}
A.~Rivero,
``Unbroken supersymmetry without new particles,''
arXiv:0910.4793 [hep-ph] (2009).

\bibitem{Rivero2024}
A.~Rivero,
``An interpretation of scalars in SO(32),''
Eur.\ Phys.\ J.\ C \textbf{84}, 1058 (2024),
arXiv:2407.05397 [hep-ph].

\bibitem{Koide1983}
Y.~Koide,
``New view of quark and lepton mass hierarchy,''
Phys.\ Rev.\ D \textbf{28}, 252 (1983).

\bibitem{Brannen2006}
C.~Brannen,
``The lepton masses,''
\url{https://brannenworks.com/MASSES2.pdf} (2006).

\bibitem{AHS2025}
S.~Antusch, K.~Hinze and S.~Saad,
``Updated Running Quark and Lepton Parameters at Various Scales,''
arXiv:2510.01312 [hep-ph] (2025).

\bibitem{Seiberg1995}
N.~Seiberg,
``Electric--magnetic duality in supersymmetric non-abelian gauge
theories,''
Nucl.\ Phys.\ B \textbf{435}, 129 (1995),
arXiv:hep-th/9411149.

\bibitem{IntriligatorSeiberg1995}
K.~Intriligator and N.~Seiberg,
``Lectures on supersymmetric gauge theories and electric-magnetic
duality,''
Nucl.\ Phys.\ Proc.\ Suppl.\ \textbf{45BC}, 1 (1996),
arXiv:hep-th/9509066.

\bibitem{RiveroZladder}
A.~Rivero,
``The $Z$-ladder and Koide-type mass relations,''
unpublished notes (2026), \texttt{related-papers/phys6\_new.tex}
in the course repository.

\bibitem{Sumino2009}
Y.~Sumino,
``Family gauge symmetry as an origin of Koide's mass formula,''
JHEP \textbf{0905}, 075 (2009).

\bibitem{HabaKoide2008}
N.~Haba and Y.~Koide,
``New origin of a bilinear mass matrix form,''
Phys.\ Lett.\ B \textbf{659}, 260 (2008),
arXiv:0708.3913 [hep-ph].

\bibitem{Ma2007}
E.~Ma,
``Lepton family symmetry and possible application to the Koide mass formula,''
Phys.\ Lett.\ B \textbf{649}, 287 (2007),
arXiv:hep-ph/0612022.

\bibitem{ISS2006}
K.~Intriligator, N.~Seiberg and D.~Shih,
``Dynamical SUSY breaking in meta-stable vacua,''
JHEP \textbf{0604}, 021 (2006),
arXiv:hep-th/0602239.

\bibitem{Rivero2013}
A.~Rivero,
``A new Koide tuple: strange-charm-bottom,''
arXiv:1111.7232 [hep-ph] (2011).

\bibitem{Zenczykowski2012}
P.~Zenczykowski,
``Remark on Koide's $Z_3$-symmetric parametrization of
quark masses,''
Phys.\ Rev.\ D \textbf{86}, 117303 (2012),
arXiv:1210.4125 [hep-ph].

\end{thebibliography}

\end{document}
